\documentclass[12pt]{article}

\usepackage{geometry}
\usepackage{titling}
\usepackage{amsmath,amssymb}
\usepackage{fontspec}
\usepackage{unicode-math}
\usepackage{array}
\usepackage[dvipsnames]{xcolor}
\usepackage{booktabs}
\usepackage{enumitem}
\usepackage{listings}

%\geometry{a4paper, top=2.5cm, bottom=2.5cm, left=2.6cm, right=2.5cm}
%\pretitle{\begin{center}\LARGE\bfseries}
%\posttitle{\end{center}\vspace{-1.5em}}
%\preauthor{\begin{center}\large}
%\postauthor{\end{center}\vspace{-1.0em}}
%\predate{\begin{center}\vspace{-0.8em}\large}
%\postdate{\par\end{center}\vspace{-0.5em}}
\setlength{\droptitle}{-2.8cm}

\setmonofont{DejaVu Sans Mono}
\setmainfont{Libertinus Serif}
\setsansfont{Libertinus Sans}
\setmathfont{STIXTwoMath-Regular.otf}

\newcommand{\Matrix}[1]{\textbf{#1}}
\newcommand{\Vector}[1]{\textbf{#1}}
\newcommand{\MMat}[1]{\mathbf{#1}}
\newcommand{\MVec}[1]{\mathbf{#1}}
\newcommand{\Norm}[1]{\lVert {\mathbf{#1}} \rVert}
\newcommand{\NormBeta}[1]{\lVert \beta \mathbf{#1} \rVert}
\newcommand{\NormConst}[1]{\lVert #1 \rVert}
% \newcommand{\Trans}{\text{\raisebox{0.5ex}{\scalebox{0.75}{T}}}}
\newcommand{\abs}[1]{\lvert {#1} \rvert}
\newcommand{\Avg}[1]{\mathbf{avg}(\MVec{#1})}
\newcommand{\Std}[1]{\mathbf{std}(\MVec{#1})}

\definecolor{lightyellow}{rgb}{1, 1, 0.8} % Define a subtle highlight color
\definecolor{codepy}{rgb}{0.9, 0.85, 0.85} %{#F8BC65}
\definecolor{codegray}{rgb}{0.5,0.5,0.5}
\definecolor{codecoffee}{rgb}{0.75,0.6,0.3}
\definecolor{codecomment}{rgb}{0.97, 0.74, 0.4}
\definecolor{lightblue}{rgb}{0.55,0.75,0.9}
\definecolor{midgray}{rgb}{0.4, 0.4, 0.3}
\definecolor{indexwhite}{rgb}{0.9, 0.9, 0.9}
\definecolor{darkgold}{HTML}{8F3F0F}
\definecolor{amber}{RGB}{255, 191, 0}
\definecolor{brickred}{RGB}{205, 92, 92}

\lstset{
    language=Python,
    basicstyle=\ttfamily\small,
    columns=fullflexible,
    frame=single,
    breaklines=true,
    showstringspaces=false
}

\newcommand{\srccodesize}{\fontsize{9pt}{10pt}\selectfont}

\lstdefinestyle{mypystyle}{
    aboveskip=2pt,
    belowskip=2pt,
    frame=single,
    framerule=0pt, % disable the frame rules
    framextopmargin=2pt,
    framexbottommargin=2pt,
    lineskip=3pt,
    % expand the margins to pull everything inside the box
    framexleftmargin=4.4mm,     % Makes the left side of the black box wider
    xleftmargin=5.8mm,
    xrightmargin=4pt,
    commentstyle=\color{codecomment},
    keywordstyle=\color{lightblue},
    numberstyle=\srccodesize\color{white!40},
    stringstyle=\color{codecoffee},
    basicstyle=\ttfamily\srccodesize\color{codepy},
    breakatwhitespace=false,
    breaklines=true,
    % captionpos=b,
    keepspaces=true,
    numbers=left,
    numbersep=0pt,
    firstnumber=1,  % Forces reset to 1 for EVERY listing
    showspaces=false,
    showstringspaces=false,
    showtabs=false,
    tabsize=2,
    language=Python,
    backgroundcolor=\color{black!80},
    emph=[2]{__init__, calculate},
    emphstyle=[2]\color{funcColor},
}
\lstset{style=mypystyle}

\title{Manipal School of Information Sciences\\AME 5151 - Applied Linear Algebra Lab\vspace*{1em}\\Assignment 2\\}

\author{Academic Year: 2026-2027, Semester: 1}

\date{Tuesday, 15 September, 2026}

\begin{document}
\pagecolor{white!2}

\maketitle

\vspace*{-1em}\noindent{\color{black!70}{\rule{\textwidth}{0.8pt}}}

{\color{black!70}{
    \noindent This assignment maps to the lab course outcome \textbf{CO2}:\vspace*{2pt}\\
    \noindent \textbf{Implement models for real-life applications using linear algebra principles
                         and analyze the results from a practical perspective.}
}}

\vspace*{1.5em}
\noindent\hfill {\textbf{Total Points}: \textbf{50}}
\vspace*{-4pt}\\
\noindent{\color{black!70}{\rule{\textwidth}{0.8pt}}}\\

\section*{Semantic Tagging with GloVe and NumPy}

You have already built a {\color{blue}{Vector}} class using tuples for
storage and have implemented operations such as
{\color{blue}{add}}, {\color{blue}{scale}}, {\color{blue}{dot}}, {\color{blue}{norm}}, and
{\color{blue}{cosine\_similarity}}.

In this lab, you will apply the same vector operations to
word embeddings. We will use the pretrained GloVe model
glove-wiki-gigaword-50. Each word represented by this model is
associated with a vector in $\mathbb{R}^{50}$.

The main objective of this lab is to move from computing cosine
similarity between one pair of vectors to computing cosine similarities
between two collections of vectors.

Suppose that, after preprocessing, the input text contains \(n\) words
that are present in the model vocabulary. Arrange their embedding
vectors as the rows of a matrix $W \in \mathbb{R}^{n \times 50}$.

Similarly, arrange the vectors corresponding to 20 predefined tags as
the rows of a matrix $T \in \mathbb{R}^{20 \times 50}$.

After normalizing every row of \(W\) and \(T\) to have Euclidean norm
one, define

\[
    \widehat{W} \in \mathbb{R}^{n \times 50},
    \qquad
    \widehat{T} \in \mathbb{R}^{20 \times 50}.
\]

The matrix $S = \widehat{W}\widehat{T}^{\,T}$ has dimensions
$S \in \mathbb{R}^{n \times 20}$.

Its entries satisfy

\[
    S_{ij} = \widehat{W}_{i,:} \, \widehat{T}_{j,:}^{\,T} = \cos(W_{i,:},T_{j,:}).
\]

Thus, \(S_{ij}\) is the cosine similarity between the \(i\)-th word
in the processed text and the \(j\)-th tag.

This is the central linear-algebra operation in the lab.

\section*{Learning Objectives}

By the end of this lab, you should be able to:

\begin{enumerate}
    \item Represent a fixed vocabulary of 20 tags as a matrix
    \[
    T \in \mathbb{R}^{20 \times 50}.
    \]

    \item Preprocess an input text and represent its in-vocabulary words
    as a matrix
    \[
    W \in \mathbb{R}^{n \times 50}.
    \]

    \item Compute cosine similarity between every in-vocabulary text
    word and every tag.

    \item Construct the complete similarity matrix
    \[
    S \in \mathbb{R}^{n \times 20}.
    \]

    \item Compute the same similarity matrix in two ways:
    \begin{enumerate}[label=(\alph*)]
        \item using your own {\color{blue}{Vector}} class and Python loops;
        \item using NumPy matrix operations.
    \end{enumerate}

    \item Verify numerically that the two implementations compute the
    same cosine-similarity values, up to normal floating-point error.

    \item Use the similarity matrix to rank the 20 tags according to
    their relevance to the input text.
\end{enumerate}

\section*{Given Input and Required Components}
Use the pretrained embedding model glove-wiki-gigaword-50.
The model associates every word in its vocabulary with a
50-dimensional NumPy vector.

You must implement:
\begin{enumerate}
    \item a vocabulary of exactly 20 tags.
    \item a preprocessing pipeline for the input text.
    \item construction of the tag matrix \(T\).
    \item construction of the text matrix \(W\).
    \item construction of the similarity matrix using your
          {\color{blue}{Vector}} class.
    \item construction of the similarity matrix using NumPy.
    \item numerical verification of the two results.
    \item tag ranking using max pooling.
\end{enumerate}

\noindent{\color{red}{\textbf{Do not use a library function that directly computes the complete cosine-similarity matrix}}}.

\section*{Part A --- Choose 20 Tags}

Choose exactly 20 tags that you consider appropriate for describing
activities associated with a research university.

Each tag must represent an academic, research, educational,
professional, institutional, or campus-related concept.

A possible starting vocabulary is:
\emph{research},
\emph{innovation},
\emph{education},
\emph{university},
\emph{students},
\emph{faculty},
\emph{campus},
\emph{engineering},
\emph{medicine},
\emph{technology},
\emph{curriculum},
\emph{collaboration},
\emph{publication},
\emph{laboratory},
\emph{scholarship},
\emph{mentorship},
\emph{internship},
\emph{entrepreneurship},
\emph{accreditation},
\emph{alumni}.

You may use this list or construct your own list.

In your documentation, briefly justify your choice of tag vocabulary.

\subsection*{Single-word tags}

If a tag consists of a single word \(t\), its tag vector is

\[
\mathbf{v}_t = \texttt{model[t]}.
\]

Every single-word tag must be present in the model vocabulary.

For example:

\begin{lstlisting}
    if tag not in model.key_to_index:
        raise ValueError(f"Tag not in vocabulary: {tag}")
\end{lstlisting}

\subsection*{Multi-word tags}

You may use a multi-word tag such as

\begin{center}
\texttt{higher education}.
\end{center}

Suppose a tag consists of \(k\) words

\[
t_1,t_2,\ldots,t_k.
\]

Represent the tag by the arithmetic mean of the individual word
vectors:

\[
\mathbf{v}_{\mathrm{tag}}
=
\frac{1}{k}
\sum_{r=1}^{k}
\mathbf{v}_{t_r}.
\]

Every component word of a multi-word tag must be present in the model
vocabulary. If any component word is absent, choose a different tag.

Document all multi-word tags used in your vocabulary.

\section*{Part B --- Build the Input Text}

Collect at most 400 words of real text concerning Manipal University
from LinkedIn posts.

Cite the source of every excerpt used.

Save the collected material in a plain-text file with extension
\texttt{.txt}.

The text may contain:

\begin{itemize}
    \item punctuation;
    \item names;
    \item numbers;
    \item hashtags;
    \item repeated words;
    \item words that do not occur in the embedding vocabulary.
\end{itemize}

Your program must handle such input without assuming that every token
has a corresponding embedding vector.

\section*{Part C --- Preprocess the Text}

Apply the following preprocessing steps in the stated order:

\begin{enumerate}
    \item convert the text to lowercase;
    \item remove punctuation;
    \item split the resulting text into tokens;
    \item remove stopwords using a small stopword set that you define;
    \item remove tokens whose length is less than or equal to 2;
    \item determine which remaining tokens occur in the embedding
          vocabulary.
\end{enumerate}

Do not use an external stopword library. Define a small stopword set
yourself.

For example:

\begin{lstlisting}
    STOPWORDS = {
        "the", "a", "an", "and", "or", "of", "to",
        "in", "on", "for", "is", "are", "was", "were",
        "with", "at", "by", "from"
    }
\end{lstlisting}

This starter set is deliberately incomplete. Real LinkedIn text will
also contain pronouns, demonstratives, and other function words
(\texttt{it}, \texttt{this}, \texttt{that}, \texttt{our}, \texttt{we},
and similar). Expand the set as needed for your own text and note in
your documentation what you added.

Let

\[
x_1,x_2,\ldots,x_m
\]

denote the tokens remaining after the first five preprocessing steps.

Some of these tokens may not be present in the model vocabulary.

Define the in-vocabulary token sequence as

\[
w_1,w_2,\ldots,w_n,
\qquad n \leq m,
\]

where each \(w_i\) is present in the model vocabulary.

Your program must report:

\begin{itemize}
    \item the number of tokens before preprocessing;
    \item the number of tokens after preprocessing;
    \item the number of in-vocabulary tokens;
    \item the number of out-of-vocabulary tokens;
    \item the distinct out-of-vocabulary tokens.
\end{itemize}

Repeated in-vocabulary words must \emph{not} be removed. If a word
appears several times in the processed text, its vector must appear
the same number of times in the text matrix.

\newpage
\section*{Part D --- Build the Tag and Text Matrices}

Implement the following functions.

\begin{lstlisting}
    def build_tag_matrix(model, tags):
        """
        Construct the tag matrix.

        Returns:
            tag_names:
                list[str] containing exactly 20 tag names.

            T:
                NumPy array with shape (20, 50).

        For a single-word tag:
            the corresponding row of T is model[tag].

        For a multi-word tag:
            the corresponding row is the arithmetic mean
            of the vectors of all component words.

        Every word used in a tag must occur in the model
        vocabulary. Otherwise, raise an error.
        """

        def build_text_matrix(model, tokens):
        """
        Construct the text matrix.

        Parameters:
            tokens:
                preprocessed tokens obtained after lowercasing,
                punctuation removal, stopword removal, and the
                length filter.

        Returns:
            in_vocab_tokens:
                tokens that occur in the model vocabulary,
                preserving their original order and repetitions.

            oov_tokens:
                tokens that do not occur in the model vocabulary.

            W:
                NumPy array with shape (n, 50), where n is the
                number of in-vocabulary tokens.

        Row i of W is the embedding vector of
        in_vocab_tokens[i].
        """
\end{lstlisting}

If

\[
w_1,w_2,\ldots,w_n
\]

are the in-vocabulary tokens, then

\[
W =
\begin{bmatrix}
\text{---} & \mathbf{w}_1^T & \text{---}\\
\text{---} & \mathbf{w}_2^T & \text{---}\\
    & \vdots          &    \\
\text{---} & \mathbf{w}_n^T & \text{---}
\end{bmatrix}
\in \mathbb{R}^{n\times 50}.
\]

Similarly,

\[
T =
\begin{bmatrix}
\text{---} & \mathbf{t}_1^T & \text{---}\\
\text{---} & \mathbf{t}_2^T & \text{---}\\
    & \vdots          &    \\
\text{---} & \mathbf{t}_{20}^T & \text{---}
\end{bmatrix}
\in \mathbb{R}^{20\times 50}.
\]

\subsection*{Checkpoint}

Print:
\begin{lstlisting}
    print(T.shape)
    print(W.shape)
\end{lstlisting}

Verify that

\[
T.\mathrm{shape}=(20,50)
\]

and

\[
W.\mathrm{shape}=(n,50).
\]

If \(n=0\), the input contains no usable words and the program must
terminate with an explanatory error message rather than attempting
the similarity computation.

\section*{Part E --- Cosine Similarity Using Your Vector Class}

Recall that for nonzero vectors
\(\mathbf{u},\mathbf{v}\in\mathbb{R}^{50}\),

\[
\cos(\mathbf{u},\mathbf{v})
=
\frac{\mathbf{u}\cdot\mathbf{v}}
{\|\mathbf{u}\|_2\|\mathbf{v}\|_2}.
\]

First compute the complete similarity matrix using your own
\texttt{Vector} class.

Implement:

\begin{lstlisting}
    def similarity_matrix_vector_class(W, T):
        """
        Compute the complete cosine-similarity matrix using
        the previously implemented Vector class.

        Every entry MUST be computed by constructing two
        Vector instances and calling their cosine_similarity
        method -- not by calling NumPy's dot/norm functions
        directly on the rows of W and T.

        Returns:
            S with shape (n_words, 20), where

            S[i, j] = Vector(W[i]).cosine_similarity(Vector(T[j])).

        This implementation may use Python loops.
        NumPy matrix multiplication must not be used to
        compute the similarities in this function.
        """
\end{lstlisting}

For every \(i\) and \(j\),

\[
S_{ij}
=
\frac{
W_{i,:}\,T_{j,:}^{\,T}
}{
\|W_{i,:}\|_2
\|T_{j,:}\|_2
}.
\]

The resulting matrix must have dimensions

\[
S_{\mathrm{Vector}}
\in
\mathbb{R}^{n\times20}.
\]

This implementation serves as the reference computation using the
vector operations that you implemented earlier in the course.

\section*{Part F --- Cosine Similarity Using NumPy}

Now compute the same matrix without using a Python-level double loop
over all text-word/tag pairs.

First normalize every row of \(W\):

\[
\widehat{W}_{i,:}
=
\frac{W_{i,:}}{\|W_{i,:}\|_2}.
\]

Similarly, normalize every row of \(T\):

\[
\widehat{T}_{j,:}
=
\frac{T_{j,:}}{\|T_{j,:}\|_2}.
\]

Then compute

\[
\boxed{
S_{\mathrm{NumPy}}
=
\widehat{W}\widehat{T}^{\,T}
}
\]

using NumPy matrix multiplication.

The matrix dimensions are

\[
\underbrace{
\widehat{W}
}_{n\times50}
\;
\underbrace{
\widehat{T}^{\,T}
}_{50\times20}
=
\underbrace{
S_{\mathrm{NumPy}}
}_{n\times20}.
\]

Implement:

\begin{lstlisting}
    def similarity_matrix_numpy(W, T):
        """
        Compute all text-word/tag cosine similarities using
        NumPy row normalization and matrix multiplication.

        Returns:
            S with shape (n_words, 20).
        """
\end{lstlisting}

A suitable implementation has the following mathematical structure:

\begin{lstlisting}
    W_norm = np.linalg.norm(W, axis=1, keepdims=True)
    T_norm = np.linalg.norm(T, axis=1, keepdims=True)

    W_hat = W / W_norm
    T_hat = T / T_norm

    S = W_hat @ T_hat.T
\end{lstlisting}

Before division, verify that the row norms are nonzero.

Do not call a library cosine-similarity function for this computation.

\section*{Part G --- Verify the Two Implementations}

The two similarity matrices must represent the same mathematical
quantities.

Before comparing them, convert \(S_{\mathrm{Vector}}\) to a NumPy
array (e.g.\ \texttt{S\_vector = np.array(S\_vector)}) if your
\texttt{Vector}-class implementation returned a nested list.

Compare them using

\begin{lstlisting}
    np.allclose(S_vector, S_numpy)
\end{lstlisting}

with an appropriate floating-point tolerance.

The result should be \texttt{True}.

Because floating-point arithmetic is approximate, do not require the
corresponding entries to be bit-for-bit identical.

Also report

\begin{lstlisting}
    np.max(np.abs(S_vector - S_numpy))
\end{lstlisting}

to show the largest absolute numerical difference between the two
implementations.

\section*{Part H --- Rank Tags Using Max Pooling}

The \(j\)-th column of \(S\) contains the similarities between tag
\(j\) and every in-vocabulary word in the text:

\[
S_{1j},S_{2j},\ldots,S_{nj}.
\]

Define the max-pool score for tag \(j\) as

\[
r_j
=
\max_{1\leq i\leq n}S_{ij}.
\]

Thus,

\[
\boxed{
r_j=\max_i S_{ij}
}
\]

is the similarity between tag \(j\) and the single text word having
the highest cosine similarity to that tag.

Implement:

\begin{lstlisting}
    def rank_tags(tag_names, S, in_vocab_tokens):
        """
        For each tag j:

            score[j] = max_i S[i, j]

        Return the tags sorted in descending order of score.

        For each tag, also return the text token that produced
        the maximum score.
        """
\end{lstlisting}

For each of the top eight tags, display:

\begin{itemize}
    \item the rank;
    \item the tag;
    \item its max-pool similarity score;
    \item the text word that produced that score.
\end{itemize}

For example, the output format may be:

\begin{verbatim}
Rank   Tag             Score      Best-matching text word
----------------------------------------------------------
1      research        0.91       research
2      innovation      0.84       innovative
3      technology      0.79       digital
...
\end{verbatim}

Your numerical values will depend on the text and the selected tag
vocabulary.

\subsection*{Interpretation of Max Pooling}

Max pooling answers the following precise question:

\begin{quote}
For this tag, which single word in the document has the highest
cosine similarity to the tag?
\end{quote}

A high max-pool score does \emph{not} imply that the corresponding
topic dominates the document. A tag can receive a high score because
of a single strongly related word.

If a text word is identical to a tag (e.g.\ the word ``research''
appears verbatim and ``research'' is one of your tags), that pair
gets a similarity of exactly 1, since cosine similarity of a vector
with itself is 1. On real text this will happen often, and such exact
matches will tend to dominate the top of your ranking ahead of tags
that are only captured through a genuinely different, semantically
related word. Note in your output which top-ranked tags were matched
by an identical word versus a different one.

Therefore, max-pool ranking measures the presence of strong local
semantic evidence for a tag; it does not measure how frequently or
how extensively the corresponding topic occurs throughout the text.

\section*{Required Unit Tests}

Write unit tests for at least the following cases.

\begin{enumerate}
    \item Verify that
    \[
    T.\mathrm{shape}=(20,50).
    \]

    \item Verify that if the text contains \(n\) in-vocabulary tokens,
    then
    \[
    W.\mathrm{shape}=(n,50).
    \]

    \item Verify that
    \[
    S.\mathrm{shape}=(n,20).
    \]

    \item Select at least one pair consisting of a text word and a tag.
    Compute its cosine similarity using your \texttt{Vector} class and
    verify that it agrees with the corresponding entry of the NumPy
    similarity matrix.

    \item Verify that
    \begin{lstlisting}
        np.allclose(S_vector, S_numpy)
    \end{lstlisting}
    returns \texttt{True}.

    \item Verify that an out-of-vocabulary text token is excluded from
    \(W\) and is reported as out-of-vocabulary.

    \item Verify that a tag vocabulary containing an out-of-vocabulary
    component causes \texttt{build\_tag\_matrix} to raise an error
    (it must not silently drop the offending tag and proceed with the
    remaining 19).

\end{enumerate}

\section*{Questions for the Lab Report}

Answer the following questions precisely.

\begin{enumerate}

    \item If there are \(n\) in-vocabulary words in the processed text,
    give the dimensions of:
    \begin{enumerate}
        \item one word vector;
        \item \(W\);
        \item \(T\);
        \item \(T^T\);
        \item \(S\).
    \end{enumerate}

    \item Explain why, after row normalization,
    \[
    \widehat{W}_{i,:}
    \widehat{T}_{j,:}^{\,T}
    \]
    is equal to the cosine similarity between the corresponding
    original vectors.

    \item Explain why
    \[
    \widehat{W}\widehat{T}^{\,T}
    \]
    computes all \(20n\) pairwise text-word/tag cosine similarities
    in one matrix multiplication.

    \item Suppose a single word in the document has very high
    similarity to the tag \texttt{medicine}, while all other words
    have low similarity to that tag. Explain how max pooling treats
    this situation.

    \item Does a high max-pool score imply that the corresponding tag
    describes the dominant topic of the document? Explain.

    \item For your own top-eight ranking, identify which tags were
    matched by a text word identical to the tag itself, and which
    were matched by a genuinely different word. What does the
    difference tell you about what cosine similarity over word
    embeddings is, and is not, capturing?

    \item If a concept is not represented by any of the 20 chosen
    tags, can the ranking algorithm introduce a new tag for that
    concept? Explain.

    \item Why are out-of-vocabulary tokens excluded from \(W\)?

    \item Explain the mathematical relationship between your
    pair-by-pair \texttt{Vector.cosine\_similarity} computation and
    the NumPy matrix multiplication used to construct \(S\).

\end{enumerate}

\section*{Submission Requirements}

Submit:

\begin{enumerate}
    \item the source text file containing at most 400 words;
    \item citations for all LinkedIn excerpts used;
    \item your Python source code;
    \item your unit tests;
    \item program output showing:
    \begin{itemize}
        \item token statistics;
        \item \(T.\mathrm{shape}\);
        \item \(W.\mathrm{shape}\);
        \item verification using \texttt{np.allclose};
        \item the maximum numerical difference between the two
              similarity matrices;
        \item the top eight ranked tags and their best-matching
              text words;
    \end{itemize}
    \item answers to the lab-report questions.
\end{enumerate}

\section*{Rubric --- 50 Points}

\begin{center}
\begin{tabular}{lr}
\toprule
\textbf{Component} & \textbf{Points}\\
\midrule
Tag vocabulary and tag-matrix construction        & 8\\
Text preprocessing and OOV handling               & 7\\
Similarity matrix using the \texttt{Vector} class & 8\\
Vectorized NumPy similarity matrix                 & 10\\
Max-pool tag ranking                               & 6\\
Unit tests and numerical verification              & 7\\
Documentation and lab-report answers               & 4\\
\midrule
\textbf{Total}                                     & \textbf{50}\\
\bottomrule
\end{tabular}
\end{center}

\end{document}
